\scanpage{25}
\mspage{-2}
l'application de $\handwritten{\otimes}^{p}\mathbb{A}$ dans $\mathbb{A}$ qui est définie par la multiplication $m$ de $A$.

\noindent - $\handwritten{\Sigma}^{p}\mathbb{A}$ est la sous-algèbre de $\handwritten{\otimes}^{p}\mathbb{A}$ qui est formée des tenseurs symétriques d'ordre $p$. Comme la diagonale $\handwritten{\Delta}$ est commutative, $\handwritten{\Delta}^{p}$ est le composé d'une application $a:\mathbb{A}\mathrel{\text{-{}-{}-{}-{}-{}-}\!\handwritten{>}}\handwritten{\Sigma}^{p}\mathbb{A}$ et de l'injection canonique $i:\handwritten{\Sigma}^{p}\mathbb{A}\mathrel{\text{-{}-{}-{}-{}-{}-}\!\handwritten{>}}\handwritten{\otimes}^{p}\mathbb{A}$.

Si $(a_i)$ est une base de $\mathbb{A}$, les produits $b=a_{i_1}\handwritten{\otimes}\ldots\handwritten{\otimes}a_{i_r}$ forment une base de $\handwritten{\otimes}^{p}\mathbb{A}$. Si $t_b$ est le sous-groupe du groupe symétrique $s_p$ formé des permutations laissant fixe $b$, nous posons $S(b)=\handwritten{\sum}\handwritten{\sigma}_{\deleted{\unsure{i}}j}(b)$, les $\handwritten{\sigma}_{\deleted{\unsure{i}}j}$ parcourant un système de représentants de $s_p\handwritten{/}t_b$. Il est clair que les $S(b)$ forment une base de $\handwritten{\Sigma}^{p}\mathbb{A}$.

\noindent - $S^{p}\mathbb{A}$ est l'image de $\handwritten{\otimes}^{p}\mathbb{A}$ dans l'algèbre symétrique de $A$. Comme la mutiplicati\unsure{\deleted{p}o|\deleted{o}p}n $m$ est comm\deleted{j}utative, $m^p$ est le composé de la surjection canonique $q:\handwritten{\otimes}^{p}\mathbb{A}\mathrel{\text{-{}-{}-{}-{}-{}-}\!\handwritten{>}}S^p\mathbb{A}$ et d'une appli\deleted{\unsure{s|x}}cation $b:S^p\mathbb{A}\mathrel{\text{-{}-{}-{}-{}-{}-{}-}\!\handwritten{>}}\mathbb{A}$.

Si $(a_i)$ est une base de $\mathbb{A}$, $S^p\mathbb{A}$ est l'ensemble des polynô mes homogènes de degré $p$ en les $a_i$.

\noindent - Nous allons définir une application $r:\handwritten{\Sigma}^{p}\mathbb{A}\mathrel{\text{-{}-{}-{}-{}-{}-}\!\handwritten{>}}\mathbb{A}^{(p)}$ en exhibant les $r(S(b))$ :

si $a_{i_1}=\ldots=a_{i_r}=a$, nous posons $r(S(b))=r(a\handwritten{\otimes}\ldots\handwritten{\otimes}a)=1\handwritten{\otimes}a$.

si les $a_{i_j}$ ne sont pas tous égaux, nous posons $r(S(b))=r(\handwritten{\sum}\handwritten{\sigma}_j(b))=0$.

\noindent - Nous définissons de même une application $j:\mathbb{A}^{(p)}\mathrel{\text{-{}-{}-{}-{}-{}-}\!\handwritten{>}}S^p\mathbb{A}$ telle que
\[
j(1\handwritten{\otimes}a_i)=a_i^p=q(a\handwritten{\otimes}\ldots\handwritten{\otimes}a)\ .
\]
% archive sheet number: 24
